% Options for packages loaded elsewhere
\PassOptionsToPackage{unicode}{hyperref}
\PassOptionsToPackage{hyphens}{url}
\documentclass[
  ignorenonframetext,
]{beamer}
\newif\ifbibliography
\usepackage{pgfpages}
\setbeamertemplate{caption}[numbered]
\setbeamertemplate{caption label separator}{: }
\setbeamercolor{caption name}{fg=normal text.fg}
\beamertemplatenavigationsymbolsempty
% remove section numbering
\setbeamertemplate{part page}{
  \centering
  \begin{beamercolorbox}[sep=16pt,center]{part title}
    \usebeamerfont{part title}\insertpart\par
  \end{beamercolorbox}
}
\setbeamertemplate{section page}{
  \centering
  \begin{beamercolorbox}[sep=12pt,center]{section title}
    \usebeamerfont{section title}\insertsection\par
  \end{beamercolorbox}
}
\setbeamertemplate{subsection page}{
  \centering
  \begin{beamercolorbox}[sep=8pt,center]{subsection title}
    \usebeamerfont{subsection title}\insertsubsection\par
  \end{beamercolorbox}
}
% Prevent slide breaks in the middle of a paragraph
\widowpenalties 1 10000
\raggedbottom
\AtBeginPart{
  \frame{\partpage}
}
\AtBeginSection{
  \ifbibliography
  \else
    \frame{\sectionpage}
  \fi
}
\AtBeginSubsection{
  \frame{\subsectionpage}
}
\usepackage{iftex}
\ifPDFTeX
  \usepackage[T1]{fontenc}
  \usepackage[utf8]{inputenc}
  \usepackage{textcomp} % provide euro and other symbols
\else % if luatex or xetex
  \usepackage{unicode-math} % this also loads fontspec
  \defaultfontfeatures{Scale=MatchLowercase}
  \defaultfontfeatures[\rmfamily]{Ligatures=TeX,Scale=1}
\fi
\usepackage{lmodern}
\ifPDFTeX\else
  % xetex/luatex font selection
\fi
% Use upquote if available, for straight quotes in verbatim environments
\IfFileExists{upquote.sty}{\usepackage{upquote}}{}
\IfFileExists{microtype.sty}{% use microtype if available
  \usepackage[]{microtype}
  \UseMicrotypeSet[protrusion]{basicmath} % disable protrusion for tt fonts
}{}
\makeatletter
\@ifundefined{KOMAClassName}{% if non-KOMA class
  \IfFileExists{parskip.sty}{%
    \usepackage{parskip}
  }{% else
    \setlength{\parindent}{0pt}
    \setlength{\parskip}{6pt plus 2pt minus 1pt}}
}{% if KOMA class
  \KOMAoptions{parskip=half}}
\makeatother
\setlength{\emergencystretch}{3em} % prevent overfull lines
\providecommand{\tightlist}{%
  \setlength{\itemsep}{0pt}\setlength{\parskip}{0pt}}
% Slide layout defaults for warning-clean scientific decks.
\usepackage{etoolbox}
\IfFileExists{xurl.sty}{\usepackage{xurl}}{}
\IfFileExists{seqsplit.sty}{\usepackage{seqsplit}}{\newcommand{\seqsplit}[1]{#1}}
\protected\def\breakseq#1{\seqsplit{#1}}
\protected\def\breaktt#1{\begingroup\ttfamily\seqsplit{#1}\endgroup}
\setlength{\emergencystretch}{6em}
\tolerance=5000
\hbadness=10000
\hfuzz=1pt
\setlength{\tabcolsep}{2pt}
\AtBeginEnvironment{longtable}{\tiny\renewcommand{\arraystretch}{0.86}\setlength{\tabcolsep}{1pt}}
\AtBeginEnvironment{tabular}{\tiny\renewcommand{\arraystretch}{0.86}\setlength{\tabcolsep}{1pt}}

% Natbib and cross-reference fallbacks — slides don't load natbib
% or cleveref, but combined-PDF manuscript prose may emit these
% commands. The fallback renders citations as a bracketed key list
% and cross-references as detokenized labels so slides don't fail on
% undefined control sequences or raw underscores. \providecommand is
% a no-op if packages load later.
\providecommand{\citep}[1]{[#1]}
\providecommand{\citet}[1]{#1}
\providecommand{\citealp}[1]{#1}
\providecommand{\citeauthor}[1]{#1}
\providecommand{\citeyear}[1]{#1}
\providecommand{\cref}[1]{\texttt{\detokenize{#1}}}
\providecommand{\Cref}[1]{\texttt{\detokenize{#1}}}
\usepackage{bookmark}
\IfFileExists{xurl.sty}{\usepackage{xurl}}{} % add URL line breaks if available
\urlstyle{same}
% fallback for those not using the hyperref driver hyperxmp:
\makeatletter
\@ifundefined{xmpquote}{\newcommand{\xmpquote}[1]{#1}}{}
\makeatother
\hypersetup{
  hidelinks,
  pdfcreator={LaTeX via pandoc}}

\author{\texorpdfstring{}{}}
\date{}

\begin{document}

\begin{frame}[fragile,allowframebreaks]{Future Directions}
\protect\phantomsection\label{future-directions}
\begin{enumerate}
\tightlist
\item
  \textbf{Supply-chain provenance}: Integration with software supply
  chain frameworks such as in-toto \breakseq{[@torresarias2019intoto]} and SLSA
  (Supply-chain Levels for Software Artifacts) \breakseq{[@openssf2023slsa]} to
  provide end-to-end attestation from source commit through build
  pipeline to published artifact. SLSA's four graduated levels of build
  integrity (from basic provenance to hermetic, reproducible builds)
  provide a natural extension ladder for the template's currently
  document-centric provenance model. The template's existing
  steganographic layer embeds document-level provenance; supply-chain
  frameworks would add build-level provenance, closing the gap between
  ``this PDF was produced by this pipeline'' and ``this pipeline was
  executed with this verified source code.''
\item
  \textbf{Decentralized provenance}: Integration with IPFS or Arweave
  for immutable publication records, extending the SHA-256-based tamper
  detection to content-addressed storage networks.
\item
  \textbf{Digital signatures}: GPG or X.509 signing integrated into the
  steganographic layer, providing cryptographic non-repudiation in
  addition to tamper detection.
\item
  \textbf{Continuous integration}: GitHub Actions workflows that execute
  the pipeline on every push, with PDF artifacts as release assets and
  automated DOI registration via Zenodo.
\item
  \textbf{Multi-language support}: Extension of the Thin Orchestrator
  pattern to R, Julia, and Rust projects, enabling polyglot research
  workflows within the Two-Layer Architecture.
\item
  \textbf{Automated FAIR4RS assessment}: Periodic self-scoring against
  FAIRsoft metrics \breakseq{[@garijo2024fairsoft]}, with quality indicators
  (executability, documentation completeness, metadata richness) tracked
  as pipeline artifacts alongside test coverage and rendering status.
\item
  \textbf{Knowledge graph integration}: Connecting pipeline outputs to
  Active Inference Knowledge Base entries for automated meta-analysis
  and cross-project citation tracking.
\item
  \textbf{Formal verification}: Static analysis tooling to enforce the
  Thin Orchestrator pattern---verifying that scripts contain no
  algorithmic logic and that \texttt{src/} modules do not import from
  \texttt{scripts/}.
\item
  \textbf{Agentic research pipelines via MCP}: The \texttt{SKILL.md}
  descriptors already define the interface contracts for each
  infrastructure module; the natural next step is to expose them as MCP
  server endpoints \breakseq{[@anthropic2024mcp]}. An MCP server wrapping
  \breaktt{infrastructure.llm} would expose \texttt{query},
  \breaktt{review\_manuscript}, and \breaktt{translate\_abstract} as
  protocol-native Tools; an MCP server wrapping
  \breaktt{infrastructure.publishing} would expose
  \breaktt{publish\_to\_zenodo} and \breaktt{generate\_citation\_bibtex}.
  A research agent could then compose these Tools to execute the full
  pipeline---environment setup → test execution → analysis → rendering →
  validation → LLM review → DOI registration---without any human in the
  loop. This closes the loop opened by Lu et al.'s AI Scientist
  \breakseq{[@lu2024aiscientist]}, which demonstrated automated hypothesis
  generation and experimental iteration but relied on ad hoc laboratory
  scaffolding. \texttt{template/}'s pipeline, fully exposed as MCP
  tools, provides that scaffolding in a reproducible, versioned, and
  certified form. Longer-term, the Agent2Agent (A2A) protocol
  \breakseq{[@google2025a2a]} enables heterogeneous specialist agents---a
  statistical analyst, a figure designer, a peer-review simulator---to
  coordinate via a shared, protocol-mediated workspace built on
  precisely the kind of modular, well-documented infrastructure that
  \texttt{template/} provides.
\item
  \textbf{Research Infrastructure as Code and software citation}: The
  DevOps IaC paradigm, applied to research, means the entire manuscript
  pipeline is a version-controlled, reproducible artifact in its own
  right. Software Heritage \breakseq{[@cosmo2020softwareheritage]} provides
  persistent SWHIDs (Software Hash Identifiers) for source-code
  snapshots, enabling \texttt{template/} itself to be cited as a
  software artifact with a DOI-equivalent stable identifier. Combining
  this with Zenodo DOI registration (already supported by
  \breaktt{infrastructure.publishing}) creates a full citation chain: the
  paper cites the data (DOI), the data provenance cites the pipeline
  (SWHID), and the pipeline cites the framework release (Zenodo DOI).
  This three-link citation chain operationalizes the Katz et al.
  \breakseq{[@katz2021software]} software citation principles at the
  infrastructure level.
\end{enumerate}
\end{frame}

\begin{frame}[fragile,allowframebreaks]{Conclusion}
\protect\phantomsection\label{conclusion}
\texttt{template/} demonstrates that high-integrity, reproducible
research need not be onerous. By embedding provenance, testing, and
documentation into the architecture itself---rather than layering them
atop a fragmented workflow---the template transforms ``best practices''
from aspirational guidelines into enforced invariants
\breakseq{[@wilson2017good; @sandve2013ten; @lamprecht2020towards]}. The
Two-Layer Architecture ensures that infrastructure improvements
propagate to all projects without coupling. The Zero-Mock policy ensures
that tests
\href{05a_zeromock_tradeoff.md\#the-zero-mock-tradeoff}{reflect
reality}. The steganographic provenance layer ensures that published
artifacts carry their own
\href{07_security_provenance.md\#security-and-provenance}{authentication}.
The \href{08f_appendix_matrix.md\#appendix-matrix}{comparative analysis}
confirms that no existing tool integrates all eleven distinctive
capabilities---testing enforcement, coverage thresholds, cryptographic
provenance, steganographic watermarking, multi-project management,
AI-agent documentation, the agentic skill protocol, interactive TUI,
Zero-Mock policy, manuscript rendering, and pipeline
orchestration---within a single enforced pipeline.

The template is not merely a build tool; it is an epistemological
commitment. It asserts that a research paper is not a static document
but a build artifact---reproducible, verifiable, and traceable to the
code that generated it. As Knuth observed, programs should be written
for humans to read and only incidentally for machines to execute
\breakseq{[@knuth1984literate]}. We extend this dictum: research manuscripts
should be \emph{built} for verification and only incidentally for
reading. In an era of generative AI, AI-native research workspaces, and
synthetic media---where the boundary between human-authored and
machine-generated text grows increasingly indeterminate
\breakseq{[@gruenpeter2021research]}---the provenance chain from source code to
published PDF is not an administrative convenience. It is the epistemic
ground on which scientific trust must be rebuilt. That this manuscript
was itself built, tested, and watermarked by the pipeline it
describes---its metrics computed from the repository it inhabits, its
figures rendered by the code it documents---is not a rhetorical device
but a structural proof: the system works because you are reading its
output.
\end{frame}

\end{document}
